% !TeX root = ../tfg.tex
% !TeX encoding = utf8
%
%*******************************************************
% Introducción
%*******************************************************

% \manualmark
% \markboth{\textsc{Introducción}}{\textsc{Introducción}} 

\chapter{Resumen}

\paragraph*{PALABRAS CLAVE:} redes neuronales, xai, {ia explicable}, clasificación, explicabilidad, inteligencia artificial, aprendizaje automático, {visión por computador}, aprendizaje profundo, {cáncer de próstata}, {gleason score}, gleason groups.

\textcolor{red}{Selecciona menos palabras clave, 5-6 como máximo, menos es más. Lo mismo, menos anglicismos. XAI e IA explicable son lo mismo por ejemplo. Te rodeo con paréntesis en el .tex mis preferidas. Son 4, si quieres añade alguna más.}

Este TFM ocupa\textcolor{red}{abarca} el problema de mejorar la interpretabilidad\textcolor{red}{explicabilidad, ya que interpretabilidad sería cambiar la estructura para que su funcionamiento fuera más simple. Nosotros usamos XAI para entender mejor un modelo complejo que no modificamos} para modelos de clasificación de lensiones cancerosas en biopsias de próstata. El cáncer de próstata es el segundo tipo de cáncer más frecuente en varones a nivel mundial. Un diagnóstico temprano de este es crucial para reducir la mortalidad del mismo. La biopsia de próstata es el procedimiento utilizado para el diagnóstico del mismo en caso de sospecha. En los últimos años el creciente uso de técnicas de Inteligencia Artifical (IA) en el campo de la salud ha collevado mejoras significativas en el diagnóstico del cáncer de próstata.

Por otro lado, la Inteligencia Artificial Explicable\textcolor{red}{eXplicable, ya que es la X la que se utiliza para el acrónimo} (XAI) es un campo de estudio que busca mejorar la interpretabiliad de los modelos de Machine Learning (ML)\textcolor{red}{aprendizaje automático}. El uso de esta es especialmente importante en campos como la medicina, donde la toma de de decisiones repercute en la vida de personas. Para este trabajo proponemos aplicar distintas técnicas XAI para mejorar la interpretabilidad del modelo.

Para la realización de este primero explicamos conceptos claves para entenderlo y estudiamos el estado del arte de los distintos aspectos relevantes a nuestra tarea. Estos son modelos de clasificación de imágenes, métodos XAI para producción de explicaciones, métricas para medir la calidad explicaciones y métodos de regularización XAI.

Para clasificación de imágenes hemos utilizado la red convolucional EfficientNet, que es el estado del arte en poder de clasificación respecto al número de parámetros del modelo. Para producción de las explicaciones de nuestro modelo hacemos uso del método LIME que produce explicaciones locales y lineales a partir de la entrada al modelo. Esto nos permite medir la calidad de las explicaciones generadas de forma cuantitativa con las métricas REVEL, que miden distintos aspectos de explicaciones locales y lineales. 

Partiendo del método X-SHIELD de regularización XAI, que representa el estado del arte, proponemos 3 métodos novedosos de regularización XAI FX-SHIELD, FR-SHIELD y H-SHIELD. Los métodos de regularización añaden una penalización al modelo durante el entrenamiento. El primero, FX-SHIELD, añade una penalización según varíe la predicción del modelo al quitar las características menos importantes extraídas por la parte convolucional del modelo. FR-SHIELD por su parte hace lo mismo que FR-SHIELD pero quitando características aleatorias sin tener en cuenta su importancia. Finalmente H-SHIELD combina FX-SHIELD y X-SHIELD para intentar combinar las mejoras introducidas por cada propuesta.

Para el entrenamiento del modelo con los distintos métodos se dispone de un conjunto de datos de 126.109 imágenes de tamaño $224 \times 224$ correspondientes a parches de biopsias de próstata. Cada una de estas imágenes está etiquetada con una de catorce posibilidades. Cada posible etiqueta representa diferentes tipo de tejido prostático según la morfología de estas y si es tejido maligno o benigno. Esta será la clase que busca predecir el modelo que entrenamos para un parche de entrada. De estas hemos utilizado 113.491 imágenes ($90 \%$) para entrenar el modelo y 12.618 ($10 \%$) para evaluar el rendimiento en clasificación de los modelos.

Los resultados obtenidos muestran que los métodos X-SHIELD y H-SHIELD obtienen la mejor \textit{accuracy} en validación y test, superando al baseline en un $0.7\%$ y $0.6\%$ respectivamente. En cambio, FX-SHIELD y FR-SHIELD consiguen mantener la \textit{accuracy} del modelo base sin utilizar regularización. Sin embargo, FX-SHIELD destaca claramente en la calidad de las explicaciones sobre los demás, logrando mejores resultados en dos de las cinco métricas REVEL y obteniendo los mismo resultados que el resto en otras dos de las cinco. Los test bayesianos utilizados confirman que FX-SHIELD supera a X-SHIELD y H-SHIELD en \textit{prescriptivity} y en \textit{conciseness} (aunque esta última no siempre con significancia), y mejora significativamente a FR-SHIELD en \textit{conciseness}. Además, todas las variantes obtienen valores similares en \textit{local fidelity} y \textit{local concordance}. Por tanto, aunque X-SHIELD logra la mejor precisión, FX-SHIELD ofrece las explicaciones de mayor calidad, lo que es crucial en aplicaciones donde la interpretabilidad es prioritaria.


\textcolor{red}{La estructura del resumen me parece un poco desorganizada, sobretodo la primera parte. Estructura propuesta para el resumen}

\begin{itemize}
    \item Primero una introducción al problema del cáncer de próstata y la importancia de un diagnóstico temprano. Énfasis en la posibilidad de uso de técnicas de IA para mejorar el diagnóstico.
    \item Párrafo sobre la necesidad de explicabilidad en modelos de IA para escenarios de alto riesgo, como es el caso de medicina, donde las decisiones afectan a vidas humanas. Motivación para usar XAI en medicina. 
    \item Párrafo sobre los conceptos clave preliminares y el estado del arte. Sin hablar específicamente del modelo que usamos, sino de las distintas áreas que utilizamos: clasificación de imágenes para disminución de tiempos de diagnóstico, métodos XAI para producción de explicaciones, métricas para medir la calidad explicaciones y métodos de regularización XAI para mejorar la calidad de las explicaciones.
    \item Párrafo de propuestas realizadas y resumen del trabajo. Aquí ya sí que hablas ligeramente del dataset utilizado muy por encima, del modelo EfficientNet como modelo de visión, LIME como generador de explicaciones, REVEL como framework de evaluación y los métodos de regularización propuestos FX-SHIELD, FR-SHIELD y H-SHIELD. No hace falta que expliques en detalle, al final para eso está el resto del TFM.
    \item Párrafo de resultados y aprendizaje. Aquí ya sí que hablas de los resultados obtenidos pero solo a nivel muy general. No es necesario que pongas números concretos, sino que hablas del 'highlight' de los resultados. Mejora en calidad, no digas exactamente en qué métricas.
\end{itemize}

\textcolor{red}{Si ves que algunos párrafos te quedan muy largos, puedes dividirlos en varios. Para ello, centrate en que estás escribiendo un relato y cada párrafo tiene que tener una idea distinta. Por ejemplo, en el párrafo 3 puede que el relato del uso de IA + XAI con los detalles de cada área se quede largo, puedes dividirlo en dos párrafos enlazando uno con el siguiente (IA reduce tiempo pero no da explicabilidad... Y en el siguiente párrafo, aquí tienes la explicabilidad necesaria con estos métodos). Incluso XAI se puede separar en XAI para generar explicaciones, XAI para medir explicaciones, XAI para mejorar explicaciones, si es que se queda muy largo.}
\thispagestyle{empty}

\endinput
